\documentclass[11pt]{article}
\usepackage[T1]{fontenc}
\usepackage{lmodern}
\usepackage[a4paper,margin=28mm]{geometry}
\usepackage{amsmath,amssymb,amsthm,mathtools}
\usepackage{booktabs,longtable,array}
\usepackage{xurl}
\usepackage[colorlinks=true,linkcolor=black,citecolor=black,urlcolor=blue]{hyperref}
\usepackage{microtype}
\newtheorem{theorem}{Theorem}[section]
\newtheorem{lemma}[theorem]{Lemma}
\newtheorem{proposition}[theorem]{Proposition}
\newcommand{\Q}{\mathbb Q}
\newcommand{\Z}{\mathbb Z}
\newcommand{\R}{\mathbb R}
\newcommand{\F}{\mathbb F}
\newcommand{\code}[1]{\texttt{\small\detokenize{#1}}}
\newcommand{\rising}[2]{(#1)_{#2}}
\DeclareMathOperator{\cont}{cont}
\title{At Least One of $\zeta(7),\zeta(9),\ldots,\zeta(19)$ Is Irrational\\[5pt]
\large A Growing-Determinant Proof with Lean Certification}
\author{Jingwei Liu\\\small The University of Western Australia}
\date{Version 1\quad\textbar\quad 6 October 2026}
\hypersetup{pdftitle={At Least One of zeta(7), zeta(9), ..., zeta(19) Is Irrational},pdfauthor={Jingwei Liu},pdfsubject={Number theory; irrationality; Lean formalization}}
\begin{document}
\maketitle
\begin{abstract}
We prove that at least one of the seven odd zeta values from
$\zeta(7)$ through $\zeta(19)$ is irrational.
A higher-derivative rational kernel gives a sequence of
$2n\times2n$ matrices whose entries are affine rational combinations of
these values. Primitive normalization converts their determinants into
integer polynomials of degree at most $2n$. We establish the global estimate
$|\det H_n/c_n|\leq \exp(-n^2/2)$ for all sufficiently large $n$.
A reduction at primes $55n+1$ proves nonvanishing under the joint-rationality
assumption, completing the integer contradiction. The proof is formalized
in Lean using Mathlib's Riemann zeta function. This article gives the
construction, the quantitative assembly, and the correspondence with the
complete machine-checked proof distributed with this version.
\end{abstract}
\noindent\textbf{Keywords.} Odd zeta values; irrationality; hypergeometric
forms; determinants; formal verification.\\
\textbf{Mathematics Subject Classification.} 11J72, 11M06, 68V15.

\section{The seven-value theorem}
For integers $s>1$, write $\zeta(s)=\sum_{k=1}^{\infty}k^{-s}$.
Our main result is the following.
\begin{theorem}\label{thm:main}
At least one of $\zeta(7),\zeta(9),\zeta(11),\zeta(13),\zeta(15),
\zeta(17),\zeta(19)$ is irrational.
\end{theorem}
The construction belongs to the higher-derivative hypergeometric framework
developed by Zudilin~\cite{zudilin}. Public proof developments of
Devincenzis~\cite{devincenzis} and CoolRmal~\cite{coolrmal} establish the
eight-value result ending at $\zeta(21)$. Theorem~\ref{thm:main} shortens
that list by one value.

The determinant, rather than an individual linear form, is the decisive
object here. Its size grows with $n$, so both the analytic estimate and
the arithmetic cost occur on the $n^2$ scale. Under joint rationality,
clearing the seven fixed denominators costs only $Q^{2n}$. The resulting
quadratic-exponential decay contradicts nonzero integrality.

We use the $158$-centered staircase kernel described below. This parameter
family and related local denominator refinements occur in the development
notes associated with~\cite{devincenzis}. The contribution of the present
proof is the growing-determinant construction and its completed arithmetic,
analytic, and nonvanishing estimates for the seven-value list. The attached
Lean development supplies the detailed proofs of the estimates used in
Sections~\ref{sec:rates} and~\ref{sec:gate}; Section~\ref{sec:formal}
identifies their exact declarations.

\section{The rational kernel and matrix}
For $m\geq0$, let $\rising{x}{m}=\prod_{j=0}^{m-1}(x+j)$. Define
\begin{align}
\kappa_n&=2\frac{\prod_{b=0}^{15}((52-2b)n)!}
 {\prod_{a=0}^{4}((48+a)n)!^2},\label{eq:kappa}\\
R_n(t)&=\kappa_n(t+79n+1)
 \frac{\displaystyle\prod_{a=0}^{4}
 \rising{t+1}{(48+a)n}\rising{t+158n+2-(48+a)n}{(48+a)n}}
 {\displaystyle\prod_{b=0}^{15}
 \rising{t+(53+b)n+1}{(52-2b)n+1}}.\label{eq:kernel}
\end{align}
Its degree at infinity is $-92n-15$. The numerator vanishes to order at
least five at every $-k$ with $1\leq k\leq48n$, and the denominator is
nonzero at those points.

Use the even, integer-valued polynomial basis
\begin{equation}\label{eq:basis}
E_0(x)=1,\qquad
E_i(x)=\frac{2x^2\prod_{j=1}^{i-1}(x^2-j^2)}{(2i)!}\quad(i\geq1).
\end{equation}
For $i\geq1$ the identity
$E_i(x)=\binom{x+i}{2i}+\binom{x+i-1}{2i}$ proves integer-valuedness on
$\Z$. Set $d=2n$ and, for $0\leq i,j<d$, put
\begin{align}
R_{n,i,j}(t)&=R_n(t)E_i(t+79n+1)E_j(t+79n+1),\\
H_n(i,j)&=\frac1{4!}\sum_{v=0}^{\infty}R_{n,i,j}^{(4)}(v-48n).
\label{eq:matrix}
\end{align}
Every entry kernel has degree at infinity at most $-84n-19$.
The fourth-derivative series in~\eqref{eq:matrix} therefore converges
absolutely. All its evaluation points lie strictly to the right of its poles.

\subsection{The seven coordinates}
Expand the entry kernel in its unique partial fractions,
\begin{equation}\label{eq:pf}
R_{n,i,j}(t)=\sum_k\sum_{\ell=1}^{m_{n,k}}
 \frac{a_{n,i,j,k,\ell}}{(t+k)^\ell},\qquad m_{n,k}\leq16,
\end{equation}
with rational coefficients. The sums over $k$ are finite and supported
between $53n+1$ and $105n+1$. Write
$h_s(K)=\sum_{v=1}^{K}v^{-s}$. Termwise differentiation and summation give
\begin{equation}\label{eq:entry-eval}
H_n(i,j)=\sum_k\sum_{\ell=1}^{m_{n,k}}
 a_{n,i,j,k,\ell}\binom{\ell+3}{4}
 \bigl(\zeta(\ell+4)-h_{\ell+4}(k-48n-1)\bigr).
\end{equation}
The reflection $t\mapsto-t-158n-2$ makes the entry kernel odd. Pairing
reflected poles cancels the even zeta coordinates. The degree at infinity
is at most $-2$, which gives
$\sum_k a_{n,i,j,k,1}=0$, cancelling the $\zeta(5)$ coordinate. Consequently
\begin{equation}\label{eq:seven}
H_n(i,j)=b_{n,i,j}+\sum_{r=0}^{6}b_{n,i,j,r}\zeta(7+2r),
\qquad b_{n,i,j},b_{n,i,j,r}\in\Q.
\end{equation}
The constant in~\eqref{eq:seven} includes the entire harmonic correction
in~\eqref{eq:entry-eval}.

Let $\mathcal H_n(\mathbf X)$ denote the matrix obtained from
\eqref{eq:seven} by replacing $\zeta(7+2r)$ with $X_r$, and set
$D_n(\mathbf X)=\det\mathcal H_n(\mathbf X)$. This is a rational polynomial
in seven variables of total degree at most $2n$. For
$\mathbf z=(\zeta(7),\ldots,\zeta(19))$, the exact identity is
\begin{equation}\label{eq:det-eval}
D_n(\mathbf z)=\det H_n.
\end{equation}

\section{Primitive normalization and the global estimate}\label{sec:rates}
When $D_n\ne0$, choose $c_n\in\Q^\times$ and
$P_n\in\Z[X_0,\ldots,X_6]$ with content one such that
\begin{equation}\label{eq:normalization}
D_n=c_nP_n,\qquad \deg P_n\leq2n.
\end{equation}
For the zero polynomial use $c_n=1$ and $P_n=0$. Thus $c_n\ne0$ for every
$n$, and
\begin{equation}\label{eq:same-object}
P_n(\mathbf z)=\frac{\det H_n}{c_n}
\end{equation}
holds for every index. Define the logarithmic arithmetic cost
$C_n=-\log|c_n|$ and the constant
\begin{equation}\label{eq:analytic-constant}
A=\frac{8123483}{8000}-12+8\log4.
\end{equation}

\begin{proposition}[Certified arithmetic and analytic estimates]\label{prop:rates}
For every $\epsilon>0$, the following inequalities hold for all sufficiently
large $n$:
\begin{align}
C_n&\leq(1012+\epsilon)n^2&&\text{if }D_n\ne0,\label{eq:arith}\\
|\det H_n|&\leq\exp\bigl((-A+2+\epsilon)n^2\bigr).
&&\label{eq:analytic}
\end{align}
Moreover $A>1014+\tfrac12$.
\end{proposition}

The arithmetic estimate uses the prime-by-prime Gauss valuations of the
same polynomial $D_n$. Its three certified contributions have upper
quadratic rates $74$, $433$, and $505$, which sum to $1012$. The assembly
retains exceptional primes and endpoint corrections. The passage from
prime sums to rates uses the bundled formal proof of the prime number theorem.

For the analytic estimate, a change to a scaled monic even basis preserves
the determinant up to its explicit leading-coefficient normalization.
Bounds on modulus moments give a product estimate over the $2n$ rows.
The compact phase bounds, the infinite tails, and finite-index drift are
included in that estimate. A uniform allowance of one per row contributes
the $+2$ on the quadratic scale. The determinant's factorial prefactor and
polynomial prefactors contribute $o(n^2)$; the exact basis normalization
contributes the remaining terms in~\eqref{eq:analytic-constant}.

These statements are proved by the declarations
\code{h158ArithmeticCertifiedRate} and
\code{h158AnalyticCertifiedRate}. Their source dependencies contain the
finite rational certificates and their all-index assembly. The numerical
constant comparison is the theorem
\code{h158ConservativeBudget_slack_gap}.

\begin{theorem}[Global determinant estimate]\label{thm:global}
For all sufficiently large integers $n$,
\begin{equation}\label{eq:global}
\left|\frac{\det H_n}{c_n}\right|
 =|P_n(\mathbf z)|\leq e^{-n^2/2}.
\end{equation}
\end{theorem}
\begin{proof}
If $D_n=0$, identity~\eqref{eq:det-eval} makes the left-hand side zero.
Otherwise apply Proposition~\ref{prop:rates} with
$\epsilon=(A-1014-\tfrac12)/3>0$. Since $|c_n|^{-1}=e^{C_n}$,
\[
\left|\frac{\det H_n}{c_n}\right|
\leq\exp\bigl((1014-A+2\epsilon)n^2\bigr)
\leq\exp(-n^2/2).
\]
Equation~\eqref{eq:same-object} supplies the polynomial formulation.
\end{proof}

\section{Nonvanishing at auxiliary primes}\label{sec:gate}
The next argument supplies nonvanishing for rational specializations,
precisely as needed in the proof by contradiction.
\begin{proposition}[Auxiliary-prime nonvanishing]\label{prop:gate}
Let $p=55n+1$ be prime, and let $\mathbf x\in\Q^7$ have coordinates
integral at $p$. Then
\[
D_n(\mathbf x)\ne0,\qquad v_p(D_n(\mathbf x))=-4n.
\]
\end{proposition}
The entrywise local computation proves that $p^2\mathcal H_n(\mathbf x)$
is integral at $p$ and that its reduction satisfies
\begin{equation}\label{eq:gate-reduction}
\overline{p^2\mathcal H_n(\mathbf x)}=V_nW_nV_n^{\mathsf T}
\quad\text{over }\F_p,
\end{equation}
where $W_n$ is diagonal with every diagonal entry nonzero. The entries of
$V_n$ are the reductions of $E_i(u_j)$ at
\[
u_j=24n+1+j,\qquad 0\leq j<2n.
\]
The complete local reduction is established in
\code{H158GateAssembly.lean}; its weights are defined there from the
Laurent coefficients of the concrete kernel~\eqref{eq:kernel}.

To see why this reduction proves the proposition, observe that
$0<u_j<p$ and $u_j+u_k<p$. Hence the squares $u_j^2$ are pairwise distinct
modulo $p$. Each $E_i(u)$ is a degree-$i$ polynomial in $u^2$ with nonzero
leading coefficient modulo $p$, because $2i<p$. Thus $V_n$ is an
invertible change of basis of a Vandermonde matrix. Taking determinants
in~\eqref{eq:gate-reduction} shows that $p^2\mathcal H_n(\mathbf x)$
has unit determinant in the local ring. As its size is $2n$,
\[
0=v_p\!\left(\det(p^2\mathcal H_n(\mathbf x))\right)
 =4n+v_p(D_n(\mathbf x)),
\]
which proves both assertions.

\section{Proof of Theorem~\ref{thm:main}}
Suppose that all seven coordinates of $\mathbf z$ are rational. Choose a
positive integer $Q$ such that $Qz_r\in\Z$ for every $0\leq r\leq6$.
By the degree bound in~\eqref{eq:normalization},
\begin{equation}\label{eq:integer}
I_n=Q^{2n}P_n(\mathbf z)\in\Z
\end{equation}
for every $n$. This follows monomial by monomial: a monomial of degree
$k\leq2n$ acquires an integral coefficient after multiplication by
$Q^{2n}$.

Dirichlet's theorem on primes in arithmetic progressions gives infinitely
many primes $p\equiv1\pmod{55}$. Omitting the finitely many prime divisors
of $Q$ leaves an unbounded set of indices $n$ with $p=55n+1$ prime and
$p\nmid Q$. At each such index the coordinates of $\mathbf z$ are integral
at $p$, and Proposition~\ref{prop:gate} gives $D_n(\mathbf z)\ne0$.
Equations~\eqref{eq:normalization} and~\eqref{eq:integer} therefore give
$I_n\ne0$.

On this same unbounded set, Theorem~\ref{thm:global} implies
\[
0<|I_n|\leq\exp\bigl(2n\log Q-n^2/2\bigr)<1
\]
for all sufficiently large $n$. This is impossible for an integer.
Theorem~\ref{thm:main} follows.\hfill$\square$

\section{Lean certification and reproducibility}\label{sec:formal}
The final declaration is
\begin{quote}\small
\code{OddZetaMixed.exists_irrational_riemannZeta_seven_to_nineteen}.
\end{quote}
It asserts the existence of an exponent in $\{7,9,11,13,15,17,19\}$ for
which the real part of Mathlib's \code{riemannZeta} is irrational. At these
real arguments the complex Riemann zeta function equals its real Dirichlet
series. The declaration has no additional mathematical hypotheses. Its
transitive axiom dependencies are \code{propext}, \code{Classical.choice},
and \code{Quot.sound}.

\begin{center}\small
\begin{tabular}{@{}p{0.28\textwidth}p{0.65\textwidth}@{}}
\toprule
Proof component & Source in \code{OddZetaMixed/}\\
\midrule
Kernel and basis & \code{H158Kernel.lean}, \code{EvenBasis.lean}, \code{KernelBasis.lean}\\[3pt]
Convergent matrix entries & \code{H158Functional.lean}\\[3pt]
Seven-coordinate identity & \code{H158SevenIdentity.lean}\\[3pt]
Primitive normalization & \code{H158Normalization.lean}\\[3pt]
Arithmetic rate & \code{H158ArithmeticCertified.lean}\\[3pt]
Analytic rate & \code{H158AnalyticCertified.lean}\\[3pt]
Auxiliary-prime reduction & \code{H158GateAssembly.lean}, \code{GateVandermonde.lean}\\[3pt]
Global bound and theorem & \code{H158SevenIrrational.lean}\\
\bottomrule
\end{tabular}
\end{center}
The accepted source replay recorded on 6 October 2026 passed all 300 checks,
rebuilt 256 project modules, and recorded hashes for 326 proof, configuration,
and verifier inputs. The bundled prime-number-theorem replay rebuilt 22
dependency modules and two audit files and checked 72 targets. These
receipts accompany the source in \code{audit/2026-10-06/}.

The pinned toolchain is Lean \code{v4.34.0-rc2}, with Mathlib revision
\begin{center}\small\texttt{dc4b8d60d5edb3c493c3662126b1b7ccae7d67cf}.\end{center}
The prime number theorem dependency is selected from the repository
\begin{center}\small\texttt{AlexKontorovich/PrimeNumberTheoremAnd},\end{center}
at revision
\begin{center}\small\texttt{c39a751132c88b6e8080b74c74023fd95b3d8be0}.\end{center}
The repository records the source selection and compatibility transformations.
The verifier recompiles project and bundled dependency sources while reusing
the pinned foundational library artifacts.

From a checkout with the stated toolchain, run:
\begin{quote}\ttfamily\small
lake exe cache get\\
python3 tools/check\_snapshot.py\\
python3 -m unittest test\_verify\\
python3 verify.py
\end{quote}
The snapshot check establishes byte identity with the recorded proof inputs;
the full verifier performs the Lean source replay. The formal proof and
certificate data are part of this version's research artifact.

\clearpage
\section*{Availability and acknowledgments}
The manuscript source, Lean development, dependency provenance, and audit
receipts are distributed at
\url{https://github.com/AlexxxxxLiu/odd-zeta-seven-lean}.
This version is identified by the release tag \code{v1.0.0} and its
accompanying SHA-256 inventory.
The archival DOI is \href{https://doi.org/10.5281/zenodo.23191331}
{10.5281/zenodo.23191331}. The manuscript is licensed under CC BY 4.0;
original code is licensed under Apache-2.0, with upstream notices retained.

The author acknowledges the higher-derivative constructions of Zudilin,
the public developments of Devincenzis and CoolRmal, and the Lean, Mathlib,
PrimeNumberTheoremAnd, and LeanArchitect projects. OpenAI Codex assisted
with mathematical exploration, implementation, verification, and preparation
of the manuscript. Authorship and responsibility for this work rest with
Jingwei Liu.

\begin{thebibliography}{9}
\bibitem{zudilin}
Wadim Zudilin, \emph{Arithmetic of linear forms involving odd zeta values},
Journal de Th\'eorie des Nombres de Bordeaux \textbf{16} (2004), no.~1.
\url{https://arxiv.org/abs/math/0206176}.

\bibitem{devincenzis}
Gennaro Marco Devincenzis, \emph{At least one of
$\zeta(7),\zeta(9),\ldots,\zeta(21)$ is irrational}, public proof and
Lean development, September 2026.
\url{https://github.com/gmDevi/zeta-7-21-lean}.
Palomar record \code{PALOMAR-2026-09-30-000007}, version 1.

\bibitem{coolrmal}
CoolRmal, \emph{Irrationality of at least one of
$\zeta(7),\ldots,\zeta(21)$ and of at least one of
$\zeta(9),\ldots,\zeta(33)$: A complete proof via Zudilin's
higher-derivative construction}, public proof manuscript dated
28 September 2026, with Lean development.
\url{https://github.com/CoolRmal/odd-zeta-irrationality}.

\bibitem{pnt}
PrimeNumberTheoremAnd contributors, \emph{PrimeNumberTheoremAnd},
Lean formalization.
\url{https://github.com/AlexKontorovich/PrimeNumberTheoremAnd}.

\bibitem{mathlib}
The Mathlib Community, \emph{Mathlib 4}, Lean mathematics library.
\url{https://github.com/leanprover-community/mathlib4}.
\end{thebibliography}
\end{document}
